\begin{table}[htbp]
\centering
\caption{Reproduction of published values with the CorOncoEndpoints package.}
\label{tab:reproduction}
\begin{threeparttable}
\scriptsize
\begin{tabular}{p{2.2cm}p{4.8cm}rrl}
\toprule
Source & Quantity & Published & Computed & Judgment \\
\midrule
Fleischer et al. (2009) & Corr(PFS, OS), lambda = (2, 0.5, 1.5), Section 3 &    0.52 & 0.5222 & PASS \\
Fleischer et al. (2009) & Example 1: Corr(PFS, OS) &    0.34 & 0.3425 & PASS \\
Fleischer et al. (2009) & Example 1: proportion dying without progression (simulated, n = 1e6) &   0.209 & 0.2092 & PASS \\
Fleischer et al. (2009) & Example 3: median PFS, control &     1.7 & 1.7014 & PASS \\
Fleischer et al. (2009) & Example 3: median PFS, treatment &     3.4 & 3.4011 & PASS \\
Fleischer et al. (2009) & Example 3: median OS, control &     9.2 & 9.2072 & PASS \\
Fleischer et al. (2009) & Example 3: median OS, treatment &     9.3 & 9.3044 & PASS \\
Fleischer et al. (2009) & Example 3: 25th percentile of OS, control &       4 & 4.0034 & PASS \\
Fleischer et al. (2009) & Example 3: 25th percentile of OS, treatment &     4.1 & 4.1034 & PASS \\
Fleischer et al. (2009) & Example 3: median OS, control with lambda3 = 0.0876 &     8.6 & 8.6318 & PASS \\
Fleischer et al. (2009) & Example 2: expected OS events at 47 months, exponential OS &    1147 & 1146.0563 & PASS \\
Fleischer et al. (2009) & Example 2: expected OS events at 47 months, model with death proportion 1/3 &    1201 & 1198.2594 & EXPLAINED\tnote{b} \\
Fleischer et al. (2009) & Example 2: power of the log-rank test at 1201 events (88.7\%) &   0.887 & -- & INFO\tnote{b} \\
TrialSimulator (CorrelatedPfsAndOs3 example) & Corr(PFS, OS), h01 = 0.1, h02 = 0.05, h12 = 0.12 (simulated, 1e6) &    0.65 & 0.6470 & PASS \\
TrialSimulator (CorrelatedPfsAndOs3 example) & median PFS &    4.62 & 4.6210 & PASS \\
TrialSimulator (CorrelatedPfsAndOs3 example) & median OS &    9.61 & 9.6119 & PASS \\
Erdmann et al. (2025) & Table 1, scenario 1: PFS hazard ratio &    0.72 & 0.7200 & PASS \\
Erdmann et al. (2025) & Table 1, scenario 2: PFS hazard ratio &   0.725 & 0.7250 & PASS \\
Erdmann et al. (2025) & Table 1, scenario 3: PFS hazard ratio &   0.764 & 0.7636 & PASS \\
Erdmann et al. (2025) & Table 1, scenario 4: PFS hazard ratio &     0.8 & 0.8000 & PASS \\
\bottomrule
\end{tabular}
\begin{tablenotes}
\footnotesize
\item PASS: 18, EXPLAINED: 1, FAIL: 0, INFO: 1. Tolerances are given in \texttt{inst/reproduce/output/reproduce\_summary.md} of the package.
\item [a] The paper gives the accrual as 45 patients per month with a total accrual time of approx. 29 months. With exactly 1300 / 45 months the expected number is 1198.3; with 29 months and 1305 patients it is 1202.1. The published 1201 lies between, and the exponential case above differs in the same direction.
\item [b] The paper does not state how the power was computed. The Schoenfeld formula with the hazard ratio 9/10.6 gives 0.809 at 1201 events.
\end{tablenotes}
\end{threeparttable}
\end{table}
